% Generated draft: requires mathematical and source review.
\documentclass[11pt]{article}
\usepackage[margin=25mm]{geometry}
\usepackage{fontspec}
\setmainfont{Latin Modern Roman}
\usepackage{amsmath,amssymb,mathtools,unicode-math}
\setmathfont{Latin Modern Math}
\usepackage{enumitem,needspace,xurl,hyperref}
\setlength{\parindent}{0pt}
\setlength{\parskip}{5pt}
\setlength{\emergencystretch}{3em}
\newcommand{\N}{\mathbb{N}}
\newcommand{\Z}{\mathbb{Z}}
\newcommand{\Q}{\mathbb{Q}}
\newcommand{\R}{\mathbb{R}}
\newcommand{\C}{\mathbb{C}}
\newcommand{\F}{\mathbb{F}}
\newcommand{\A}{\mathbb{A}}
\newcommand{\PP}{\mathbb{P}}
\newcommand{\E}{\mathbb{E}}
\newcommand{\eps}{\varepsilon}
\newcommand{\dd}{\,\mathrm{d}}
\begin{document}
\section*{UnsolvedMath ID: 161}
\textbf{Source problem code:} catalogue entry
\subsection{Definitions and mathematical statement}
For a finite nonempty set $B\subset\Z$, define its difference set and difference ratio by
\[
 B-B=\{b_1-b_2:b_1,b_2\in B\},\qquad D(B)=\frac{|B-B|}{|B|}.
\]
Thus $D$ denotes a difference ratio, despite the registry title's use of ``doubling.'' The supplied research request is: given a finite nonempty $A\subset\Z$, a bound $K\ge1$ with $D(A)\le K$, and an error $\varepsilon>0$, find a large structured nonempty subset $A'\subseteq A$ for which
\[
 D(A')\le K+\varepsilon
\]
is ``obvious.'' The finite nonempty setting is needed for these cardinality ratios; the bare source statement does not explicitly give that hypothesis. Neither a required lower bound on $|A'|$, a class of permitted structures, nor a meaning of ``obvious'' is specified. Consequently this is a research programme, not a uniquely quantified conjecture. A precise quantitative target would have to specify those choices, for example a positive function $f(K,\varepsilon)$ in a requirement $|A'|\ge f(K,\varepsilon)|A|$, together with the structural class and its certificate. No particular such function or class is asserted here as the intended source theorem.

The supplied literature triage records the following concrete special case of Eberhard--Green--Manners. For each $\eta>0$, there is $c_\eta>0$ such that if a finite nonempty $A\subset\Z$ satisfies
\[
 |A-A|\le(4-\eta)|A|,
\]
then there is a finite arithmetic progression
\[
 P=\{u+jd:0\le j<L\},\qquad u,d\in\Z,\ d\ne0,\ L\ge1,
\]
with $|P|=L\ge c_\eta|A|$ and $|A\cap P|>|P|/2$. The corresponding subset $A'=A\cap P$ has positive proportional size and the transparent bound $D(A')\le4$. This special-case bound is not claimed to give $D(A')\le K+\varepsilon$ for every smaller $K$. The source describes the size constant as essentially ineffective. Extending a suitably specified version to a stated $K$ range, with effective dependence if desired, remains the programme described by the triage.

The registry's main status field says open; its literature assessment dated 17 August 2026 says partially solved, because of this special case and the lack of a determinate general target. It also explicitly states that the background's algorithmic Polynomial Freiman--Ruzsa gloss is unsupported by the source statement.
\subsection{Short English statement}
For an integer set with bounded difference ratio, seek a large structured subset whose nearly unchanged ratio is easy to certify. The source leaves ``large,'' ``structured,'' and ``obvious'' undefined, while recording an essentially ineffective progression-density theorem below difference ratio $4$.
\subsection{Sources}
\begin{itemize}
\item Local UnsolvedMath registry, permanent ID $161$, source code \texttt{catalogue entry}, ``Structured Subsets with Bounded Doubling,'' including its supplied literature triage checked 17 August 2026.
\item Sean Eberhard, Ben Green and Freddie Manners, \emph{Sets of integers with no large sum-free subset}, Annals of Mathematics \textbf{180} (2014), 621--652, Section 6; arXiv:1301.4579, \url{https://arxiv.org/abs/1301.4579}, as cited for the progression-density theorem in the record.
\item Ben Green, \emph{100 Open Problems}, Problem 30, most recent update December 2025 according to the supplied triage, \url{https://people.maths.ox.ac.uk/greenbj/papers/open-problems.pdf}. The source problem number cited by the triage is distinct from the preserved registry code.
\end{itemize}
\end{document}
